% This file compiles with both LuaLaTeX and XeLaTeX
\documentclass[11pt]{article}

% Change "review" to "final" to generate the final (sometimes called camera-ready) version.
% Change to "preprint" to generate a non-anonymous version with page numbers.
\usepackage[review]{acl}


\usepackage{booktabs}
\usepackage{multirow}
\usepackage{graphicx}
\usepackage{tcolorbox}
\usepackage{listings}
\tcbuselibrary{breakable,skins}
\usepackage{xcolor}
% This is not strictly necessary, and may be commented out,
% but it will improve the layout of the manuscript,
% and will typically save some space.
 \usepackage{microtype}

% If the title and author information does not fit in the area allocated, uncomment the following
%
%\setlength\titlebox{<dim>}
%
% and set <dim> to something 5cm or larger.

% These font selection commands work with
% LuaLaTeX and XeLaTeX, but not pdfLaTeX.
\usepackage[english,bidi=default]{babel} % English as the main language.
\babelfont{rm}{TeXGyreTermesX} % similar to Times
%%% include whatever languages you need below this line
\babelprovide[import]{hindi}
\babelfont[*devanagari]{rm}{Lohit Devanagari}
\babelprovide[import]{arabic}
\babelfont[*arabic]{rm}{Noto Sans Arabic}


%\usepackage{polyglossia}
%\setdefaultlanguage{english}
%\setotherlanguages{arabic,russian,thai,hindi,kannada}

%%%%%

\title{CARMA: Large-Scale Arabic Reddit Corpus for Mental Health Analysis}

% Author information can be set in various styles:
% For several authors from the same institution:
% \author{Author 1 \and ... \and Author n \\
%         Address line \\ ... \\ Address line}
% if the names do not fit well on one line use
%         Author 1 \\ {\bf Author 2} \\ ... \\ {\bf Author n} \\
% For authors from different institutions:
% \author{Author 1 \\ Address line \\  ... \\ Address line
%         \And  ... \And
%         Author n \\ Address line \\ ... \\ Address line}
% To start a seperate ``row'' of authors use \AND, as in
% \author{Author 1 \\ Address line \\  ... \\ Address line
%         \AND
%         Author 2 \\ Address line \\ ... \\ Address line \And
%         Author 3 \\ Address line \\ ... \\ Address line}

\author{First Author \\
  Affiliation / Address line 1 \\
  Affiliation / Address line 2 \\
  Affiliation / Address line 3 \\
  \texttt{email@domain} \\\And
  Second Author \\
  Affiliation / Address line 1 \\
  Affiliation / Address line 2 \\
  Affiliation / Address line 3 \\
  \texttt{email@domain} \\}

\begin{document}
\maketitle
\begin{abstract}
Mental health has emerged as a critical global issue, yet Arabic NLP resources for mental health research remain scarce. We introduce \textbf{CARMA}, a large-scale corpus from Reddit (3,080 users, \textbf{\textasciitilde1M} posts), covering users who have self-reported diagnoses of mental health conditions. The high quality corpus targets 21 conditions spanning ten diagnostic categories, of which 18 yielded sufficient confirmed users for analysis. We collect data using a high-precision, LLM-powered extraction pipeline validated through manual human annotation that combines self-disclosure phrases with condition-specific keywords across four Arabic dialect groups (Gulf, Egyptian, Levantine, and Modern Standard Arabic) to identify posts containing genuine diagnostic self-reports. We employ an author-level approach that retrieves complete posting histories to enable longitudinal analysis. We benchmark the performance of current detection models and leverage Large Language Models to establish baselines for Arabic mental health classification. Our contribution is also methodological: we enhance automatic self-reported diagnosis by including human and LLM annotations as independent layers to improve quality, ensure data integrity, and reduce false positives.\end{abstract}



\section{Introduction}
\label{sec:introduction}
Mental health has emerged as a critical global challenge, with over one billion
people worldwide estimated to live with a mental disorder
\cite{who2025mental}. In recent years, social media platforms have become a valuable source of
large-scale, naturalistic text for studying mental health-related language.
Computational approaches have demonstrated that linguistic and behavioral signals
in user-generated content can serve as proxies for detecting mental health
conditions \cite{de2013predicting, chancellor2020methods}. The CLPsych workshop
series further institutionalized this line of research by establishing datasets, shared tasks and standardized evaluation frameworks \cite{coppersmith2015clpsych}.
A central methodological development is the \textit{self-reported diagnosis}
paradigm, which identifies users who publicly disclose having received a clinical
diagnosis. First introduced by \citet{coppersmith2014quantifying} on Twitter and
later extended to Reddit by \citet{cohan-etal-2018-smhd}, this approach has enabled the
construction of large-scale datasets without requiring manual clinical annotation.

Despite this progress, the overwhelming majority of existing resources target English-language data \cite{harrigian2021state}. Arabic remains substantially underrepresented in computational mental health research, despite being spoken by over 400 million people across diverse linguistic and cultural contexts. Arabic poses particular challenges for natural language processing, including rich morphology and extensive dialectal variation across regional varieties such as Gulf, Egyptian, and Levantine Arabic \cite{darwish2021panorama,habash2010introduction}. Existing Arabic mental health datasets are limited in scale, typically target only one or two conditions (most commonly depression or suicidal ideation), and rarely address dialectal variation systematically (see Table~\ref{tab:arabic_datasets}).

In this paper, we address these gaps by introducing \textbf{CARMA}, a large-scale and carefully-curated Arabic mental
health corpus collected from Reddit. The corpus follows the self-reported
diagnosis paradigm: we develop a dialect-aware, pattern-based extraction pipeline
that identifies Arabic diagnosis disclosures across four major dialect groups (Egyptian, MSA, KSA, Levantine).
Candidate disclosures are then validated through an \textit{LLM-as-Judge}
annotation framework \cite{zheng2023judging} employing three independent language
models (Qwen3, Jais2, GPT4), mitigating the risk of individual model bias and improving annotation
reliability in this low-resource setting \cite{fu-liu-2025-reliable}. We adopt
an author-level collection strategy, retrieving the complete posting histories of
identified users to support longitudinal analysis of linguistic patterns beyond
the initial disclosure posts.

The corpus targets 21 conditions spanning ten diagnostic
categories (Table~\ref{tab:arabic_datasets}).

%Three conditions, Avoidant Personality Disorder, Narcissistic Personality Disorder, and Schizotypal Personality Disorder, yielded insufficient confirmed users and are excluded from downstream analysis, leaving 18 conditions for evaluation.

Our contributions are as follows:
\begin{enumerate}
    \item We introduce \textbf{CARMA}, a large-scale Arabic mental health corpus
    collected from Reddit, targeting 21 conditions across ten diagnostic categories.
    %of which 18 yield sufficient data for analysis.

    \item We develop the first high precision, dialect aware, automatic extraction pipeline powered by LLMs to identify self disclosure in Arabic.

    \item We validate the automatic annotation pipeline using two native Arabic-speaking
    annotators and an LLM-as-Judge framework with three large language
    models.

    \item We provide a comprehensive analysis of corpus characteristics and
    baseline model performance, offering insights into the challenges of mental
    health detection in Arabic social media.

\end{enumerate}



\section{Related Work}
\label{sec:related_work}

\subsection{Mental Health Detection from Social Media}
\label{subsec:related_social_media}

Research on detecting mental health conditions from social media has grown
substantially over the past decade. Early work by \citet{de2013predicting}
demonstrated that linguistic and behavioral signals on Twitter can be used to
identify users experiencing depression. Since then, numerous studies have explored
computational methods for identifying mental health conditions from social media
language across multiple platforms and conditions.

The CLPsych workshop series \cite{coppersmith2015clpsych} introduced standardized
shared tasks and evaluation frameworks for detecting conditions such as depression
and post-traumatic stress disorder, enabling systematic comparison of computational
approaches and helping establish methodological foundations for the field.

A key development in this area is the \textit{self-reported diagnosis} paradigm,
which identifies users who publicly disclose receiving a clinical diagnosis through
pattern matching. \citet{coppersmith2014quantifying} first introduced this approach
by identifying Twitter users who disclosed mental health diagnoses via regular
expression matching. \citet{cohan-etal-2018-smhd} later extended this methodology to
Reddit, introducing the SMHD dataset, a large-scale corpus covering nine mental
health conditions. The paradigm has since been applied to other languages,
including German with the SMHD-GER dataset \cite{zanwar2023smhd}.
\citet{harrigian2021state} provide a comprehensive analysis of existing social
media datasets used in mental health research, highlighting methodological
challenges in dataset construction and emphasizing that the majority of existing
resources target English, leaving other languages substantially underrepresented.

\subsection{Arabic NLP for Mental Health}
\label{subsec:related_arabic}

Compared to English, computational research on mental health in Arabic remains
relatively limited. Arabic presents several challenges for natural language
processing, including rich morphology, substantial dialectal variation across
regional varieties such as Gulf, Egyptian, and Levantine Arabic, and limited
availability of large annotated datasets \cite{darwish2021panoramic, alayba2025arabic}.
Table~\ref{tab:arabic_datasets} summarizes existing Arabic datasets related to
mental health and situates our corpus within this landscape.

Several studies have constructed datasets from X to study specific mental health
conditions. \citet{al2024classification} introduce a dataset for detecting
obsessive-compulsive disorder symptoms in Arabic tweets. For suicidal ideation
detection, \citet{abdulsalam2024detecting} create a dataset of approximately 5.7K
tweets, while \citet{baghdadi2022optimized} provide a smaller dataset of around
2K tweets.

Depression detection has received the most attention in Arabic social media
research. \citet{hassib2022aradepsu} introduce AraDepSu, a dataset of
approximately 20K tweets annotated for depression and suicidal ideation.
Earlier work by \citet{almouzini2019detecting} developed a dataset for identifying
depressed users on Twitter. \citet{elramly2021cairodep} present CairoDep, a corpus
of approximately 7K tweets manually labeled for depression detection.
\citet{maghraby2022mood} introduce the Mood Change dataset, a Modern Standard
Arabic collection organized around PHQ-9 depression screening categories to
facilitate symptom-level analysis.

Beyond social media, \citet{alhuzali2024mentalqa} introduce MentalQA, a dataset
of 1,000 question-and-answer posts from the Altibbi medical platform annotated for
question types and answer strategies. Unlike social media datasets, MentalQA
focuses on question-answer classification rather than detecting mental health
conditions from user-generated discourse.

Despite these efforts, existing Arabic mental health datasets exhibit three
notable limitations. First, most datasets focus on only one or two conditions,
typically depression or suicidal ideation, leaving the broader spectrum of mental
health conditions largely unaddressed. Second, annotations are typically performed
at the post level, without modeling user-level context or longitudinal behavior.
Third, dialectal variation across Arabic varieties is rarely addressed
systematically during dataset construction. Our corpus is designed to address
all three of these limitations.

\subsection{LLM-as-Judge for NLP Annotation}
\label{subsec:related_llm_judge}

The \textit{LLM-as-Judge} paradigm proposes using large language models as
automated evaluators for text generation and annotation tasks.
\citet{zheng2023judging} show that strong language models can approximate human
judgments in several evaluation settings by ranking model outputs according to
task-specific criteria. As a result, LLM-based evaluation has increasingly been
adopted in NLP workflows, including dataset filtering, label verification, and
large-scale annotation pipelines.

\citet{gilardi2023chatgpt} demonstrate that ChatGPT outperforms crowd workers on
several annotation tasks, including relevance, stance, and frame detection,
suggesting that LLM-based annotation can be competitive with human annotators
on certain labeling tasks. In the mental health domain, large language models
have been applied to classification and reasoning over online mental health
discourse \cite{xu2024mental}.

However, the reliability of LLM-based judging remains an active area of
investigation. \citet{fu-liu-2025-reliable} show that agreement between LLM judges
can vary substantially across languages, with particularly low consistency in
low-resource settings. These findings directly motivate our use of three
independent LLM judges rather than a single model, and our inclusion of
validation steps to assess inter-model agreement before accepting a label.



\begin{table*}
\centering
\small
\begin{tabular}{cl l rll}
\toprule
& \textbf{Dataset} & \textbf{Condition} & \textbf{Size} & \textbf{Domain} & \textbf{Anno.} \\
\midrule
& MentalQA~\cite{alhuzali2024mentalqa} & N/A & 1,000 & Altibbi & Human \\
& OCDTweets~\cite{al2024classification} & OCD symptoms & 5K & X & Human \\
& Arabic Suicidal Tweets~\cite{abdulsalam2024detecting} & Suicidal ideation & 5.7K & X & Human \\
& Suicide Tweets~\cite{baghdadi2022optimized} & Suicidal ideation & 2.0K & X & Human \\
& AraDepSu~\cite{hassib2022aradepsu} & Depression, Suicidal ideation & 20K & X & Human \\
& Detecting Depression~\cite{almouzini2019detecting} & Depression & 2.7K & X & Human \\
& CairoDep~\cite{elramly2021cairodep} & Depression & 7K & X & Human \\
& Mood Change~\cite{maghraby2022mood} & Depression & 49K & X & Auto. \\
\midrule
& & \textit{Schizophrenia Spectrum:}& & & \\
& & Delusional Disorder, Schizophrenia & & & \\
& & \textit{Mood:} Bipolar, Depression & & & \\
& & \textit{Anxiety:} Anxiety, Panic Disorder & & & \\
& & \textit{OCD \& Related:} OCD, & & & \\
& & \textit Trichotillomania & & & \\

& & \textit{Trauma:} PTSD & & & \\
& \textbf{Ours} & \textit{Eating:} Eating Disorder & \textbf{1M} & \textbf{Reddit} & \textbf{Auto.} \\
& & \textit{Neurodevelopmental:} ADHD, Autism & & & \\
& & \textit{Dissociative:} DID & & & \\
& & \textit{Sleep:} Sleep Disorder & & & \\
& & \textit{Personality:} BPD, Narcissistic PD, & & & \\
& & ~~Avoidant PD, Schizoid PD, & & & \\
& & ~~Paranoid PD & & & \\
& & \textit{Clinical Risk:} Suicidal Ideation & & & \\
\bottomrule
\end{tabular}
\caption{Comparison of Arabic Mental Health Corpora. For our corpus, all 21 target conditions are listed; *some conditions are excluded from analysis due to insufficient confirmed users. The size here is in number of single document posts/comments.}

\label{tab:arabic_datasets}
\end{table*}



\section{Dataset Construction}
\label{subsec:data_source}

\subsection{Diagnosed Dataset}

We collect data from Reddit using pre-downloaded JSONL archives from Arctic Shift\footnote{\url{https://github.com/ArthurHeitmann/arctic_shift}} (Pushshift format), containing posts and comments from Arabic-speaking subreddits. Reddit offers longer, more detailed self-reports that are well-suited to the self-reported diagnosis paradigm. To the best of our knowledge, there are no dedicated Arabic-language mental health subreddits comparable to those in English (e.g., r/depression, r/SuicideWatch). Consequently, our dataset is derived entirely from general regional subreddits (e.g., r/saudiarabia, r/Egypt, r/jordan), from which we extract mental-health-relevant content.

We use a curated list of condition keywords (e.g., \foreignlanguage{arabic}{اكتئاب حاد}, "Severe Depression") and diagnosis phrases (e.g., \foreignlanguage{arabic}{تم تشخيصي بـ}, "I was diagnosed with"), checking whether they co-occur within a 40-character window (measured at the Unicode character level). See Appendix~\ref{app:keywords} and Appendix~\ref{app:phrases} for a detailed pattern description. To enhance inclusivity and robustness, the phrase list incorporates variations across Arabic dialects, including Modern Standard Arabic, Egyptian, and Levantine, along with frequent spellings and misspellings to account for the diverse posting habits of social media users. Condition keywords name each target condition in text, covering formal clinical terms, colloquial Arabic labels, and English terms.

\subsection{Preprocessing}

We apply two filtering steps to the collected post histories. First, we remove all posts containing mental health-related content (self-disclosures, condition keywords, and medication mentions) to prevent models from relying on explicit diagnostic markers when learning to classify users. This design encourages models to learn more subtle linguistic and behavioral signals rather than trivial keyword-based cues, aligning with prior work in mental health NLP \cite{cohan-etal-2018-smhd}. Second, we filter out non-Arabic content to ensure linguistic consistency across the corpus.

\subsection{Control Dataset}

We further construct a control dataset in a similar manner by scraping entire public user post histories from general Arabic-speaking subreddits (e.g., r/SaudiArabia) for candidate control users. Control users are subject to strict exclusion criteria: we immediately discard a candidate if they have ever posted about a mental health-related issue on the platform. For this filtration, we use the same pattern developed to identify diagnosed users, scanning for matching phrases, medication keywords, or other mental health-related terminology. See Appendix~\ref{app:phrases} for more detail. Such stringent exclusion criteria are inspired by previous literature \cite{cohan-etal-2018-smhd} to minimize the risk of including affected users in the control group.



\section{Annotation and Validation}
\label{sec:annotation}

\subsection{Human Annotation}
\label{subsec:human_annotation}

To assess the false positive rate and retrieval quality of the pipeline,
two native Arabic-speaking annotators independently labeled a stratified
sample of 100 posts from CARMA, classifying each post as \texttt{True Positive} (a genuine
self-disclosure of a mental health diagnosis) or \texttt{False Positive}
(a post that mentioned a mental health condition without constituting a
genuine diagnosis, such as through metaphor or sarcasm). The annotators agreed on 96 of 100 posts
(96.0\%), yielding Cohen's $\kappa = 0.756$, indicating substantial agreement.

\subsection{LLM-as-Judge Framework}
\label{subsec:llm_judge}

To verify whether candidate posts contain genuine first-person mental
health disclosures, we employ an LLM-as-Judge framework
\cite{zheng2023judging} using three large language models: GPT-4.1
\footnote{\url{https://openai.com/api/}} \cite{achiam2023gpt},
Qwen3-235B-A22B \cite{qwen3technicalreport}, and Jais-2-8B-Chat
\cite{jais2_2025}. Each model is prompted independently with the same
structured prompt (Appendix~\ref{app:prompt}) to classify each post as
either \texttt{<genuine>} (a real, first-person self-disclosure of the
target condition) or \texttt{<not\_confirmed>} (a false positive). The
prompt instructs each model to act as a board-certified psychiatrist
fluent in Arabic and its dialects, and to classify strictly based on
what the author explicitly states, without inferring or assuming a
diagnosis.

GPT-4.1 provides strong general-purpose multilingual reasoning
\cite{zhong2024agieval, moell2024comparing}; Qwen3-235B-A22B offers
large-scale open-weight performance with strong Arabic coverage
\cite{qwen3technicalreport}; and Jais-2-8B-Chat is a purpose-built
Arabic language model trained on Arabic and English data with a custom
Arabic-centric vocabulary, providing dialect-sensitive judgment across
Modern Standard Arabic and regional varieties \cite{jais2_2025}.
Using multiple judges mitigates individual model biases and allows us to
measure labeling consistency across systems \cite{zheng2023judging}.
A post is retained as confirmed only when all three judges unanimously
agree on \texttt{<genuine>}, prioritising precision over recall to
minimise false positives in the final corpus.

\subsection{Agreement Results}
\label{subsec:agreement}

Table~\ref{tab:iaa} reports inter-model agreement among the three LLM
judges on CARMA. GPT-4.1 and Qwen3-235B achieve the \textbf{highest}
pairwise agreement ($\kappa = 0.652$), indicating substantial alignment
between the two larger models. Jais-2-8B-Chat shows consistently lower
agreement with both peers, which we attribute to its smaller model capacity
\cite{wei2022emergent} and higher parse failure rate. Fleiss' $\kappa$
across all three models is 0.427, reflecting moderate agreement overall,
a level consistent with the inherent difficulty of distinguishing genuine
self-disclosures from figurative or sarcastic mentions in social media text.

Table~\ref{tab:iaa_human} (Appendix) compares each LLM judge against the two human
annotators. The human-to-human agreement (first row) serves as a
reference point: substantial on CARMA ($\kappa = 0.756$). The LLM Majority vote generally achieves higher or comparable
$\kappa$ with both annotators relative to any individual model, reaching
up to $\kappa = 0.590$ (Ann-2). This pattern validates the ensemble approach: aggregating
three judges mitigates the idiosyncratic errors of any single model and
brings LLM judgments closer to the human standard. Although majority
voting achieves closer alignment with human annotators, unanimous
agreement is adopted to maximize corpus precision, accepting lower
recall in favour of minimizing false positives. Full details on how pairwise $\kappa$ values and
sample sizes are computed are provided in
Appendix~\ref{app:agreement_details}.

\begin{table}[h]
\centering
\small
\begin{tabular}{lc}
\toprule
\textbf{Model Pair / Metric} & \textbf{CARMA} \\
\midrule
GPT-4.1 vs.\ Qwen3-235B       & \textbf{0.652} \\
GPT-4.1 vs.\ Jais-2-8B        & 0.299          \\
Qwen3-235B vs.\ Jais-2-8B     & 0.363          \\
\midrule
Mean pairwise $\kappa$         & 0.438          \\
Fleiss' $\kappa$ (all three)   & 0.427          \\
\midrule
Overall agreement (\%)         & 71.4\%         \\
\bottomrule
\end{tabular}
\caption{Inter-model agreement among the three LLM judges on
a stratified 100-post sample from CARMA.
$\kappa$ values follow \citet{landis1977measurement}.}
\label{tab:iaa}
\end{table}

%%%Should we select genuine posts based on majority agreement among the LLMs (at least 2 out of 3), or only when all LLMs agree??? + clear reason.
% LLM Majority votes looks like the better choice. Majority is more consistent with human judgments for both dataset!!!!
%We need to teview previous studies [LLM majority votes vs LLM Unanimous]


\section{Corpus Statistics}
\label{sec:corpus}

\subsection{Overview}
\label{subsec:corpus_overview}
The resulting corpus comprises 3,080 unique users. Although the pipeline targets 21 conditions, some conditions yielded insufficient confirmed users and are excluded from analysis. Table~\ref{tab:corpus_summary} presents the per-condition user distribution.


\begin{table}[h]
\centering
\small
\begin{tabular}{lr}
\toprule
\textbf{Condition} & \textbf{CARMA} \\
\midrule
Depression          & 1,342 \\
Anxiety             &   537 \\
ADHD                &   386 \\
OCD                 &   221 \\
Sleep Disorder      &   284 \\
Bipolar Disorder    &    66 \\
BPD                 &    79 \\
PTSD                &    66 \\
Autism              &    85 \\
Panic Disorder      &   109 \\
Suicidal Ideation   &    75 \\
Schizophrenia       &    29 \\
Eating Disorder     &    41 \\
Trichotillomania    &     9 \\
Paranoid PD         &     6 \\
DID                 &     2 \\
\midrule
\textbf{Unique users}          & \textbf{3,080} \\
\textbf{Sum across conditions} & \textbf{3,337} \\
\bottomrule
\end{tabular}
\caption{Distribution of users across 16 retained mental health
conditions in CARMA. Each cell counts unique users
confirmed by unanimous \texttt{<genuine>} agreement across all three
LLM judges. Users self-reporting multiple conditions are counted once
per condition. *Conditions with low sample size are excluded from the subsequent experiments but are still included in the official dataset
(§\ref{subsec:corpus_overview}).}
\label{tab:corpus_summary}
\end{table}


Once a user is confirmed as genuine by unanimous LLM-judge agreement, we retrieve their complete posting history to enable longitudinal analysis beyond the initial self-disclosure post.


\subsection{Language and Dialect}
Dialect identification was performed using CAMeLBERT-Mix DID NADI Model\footnote{\url{https://huggingface.co/CAMeL-Lab/bert-base-arabic-camelbert-mix-did-nadi}} \cite{inoue-etal-2021-interplay}, a BERT-based transformer fine-tuned on the NADI country-level dataset for Arabic dialect classification. As shown in Figure~\ref{fig:dialect_dist},
Egyptian Arabic is the dominant dialect in CARMA (55.75\%), with Egypt as the primary contributing country, consistent with its status as one of the most populous Arabic-speaking countries with approximately 116 million inhabitants as of 2024\footnote{Population estimates sourced from official 2024 UN demographic data.} \citep{un2024wpp_data}. Gulf Arabic constitutes the second-largest group, primarily driven by Saudi Arabia (approximately 35 million inhabitants). This distribution reflects the strong presence of Egyptian Arabic speakers on Reddit and their engagement in public regional subreddits.


\begin{figure}[h]
    \centering
    % TODO: Replace with Reddit-only dialect distribution figure
    % \includegraphics[width=\linewidth]{latex/figures/dialect_pie_reddit.png}
    \fbox{\parbox{\linewidth}{\centering\vspace{1cm}[Figure placeholder: Arabic dialect distribution in CARMA (Reddit)]\vspace{1cm}}}
    \caption{Arabic Dialect Distribution in CARMA}
    \label{fig:dialect_dist}
\end{figure}

\subsection{Comorbidity Analysis}
\label{subsec:comorbidity}

\begin{table}[t]
\centering
\begin{tabular}{lr}
\toprule
\textbf{\# Conditions} & \textbf{CARMA} \\
\midrule
1 & 3,025 \\
2 &    52 \\
3 &     2 \\
4 &     1 \\
\bottomrule
\end{tabular}
\caption{Comorbidity class counts for CARMA, adjusted to match the total unique user count (3,080 users).}
\label{tab:comorbidity_counts}
\end{table}
It is common for multiple mental health conditions to co-occur within the same individual \cite{cohan-etal-2018-smhd}. We therefore analyze comorbidity patterns across diagnosed users in our dataset. Table~\ref{tab:comorbidity_counts} reports the number of users exhibiting one or more diagnosed conditions.

The majority of users are associated with a single condition (3,025 users). A smaller subset of users exhibit comorbidity, with 52 users associated with two conditions. Higher-order comorbidities are rare, with only 2 users presenting three and 1 user presenting four co-occurring conditions.

These findings suggest that while comorbidity exists in the dataset, it remains relatively infrequent, reflecting both the natural distribution of self-reported diagnoses and the precision-oriented design of our extraction pipeline. The limited number of multi-condition users also highlights the importance of treating each condition as a distinct classification task while acknowledging potential overlap in downstream analyses.

\section{Experiments}
\label{sec:experiments}

We conducted two binary classification experiments to evaluate the utility of the dataset for downstream classification tasks. The first experiment employed classical models trained on bag-of-words features using TF-IDF representations at the post level, while the second leveraged pre-trained language models (PLMs) fine-tuned at the user level. This distinction reflects the differing capabilities of the models: classical models are well-suited for independent, high-volume post-level classification, whereas PLMs can better capture contextual and stylistic patterns when provided with richer user-level signals.

For classical models, each post is treated as an independent instance, and labels are inherited from the corresponding user. For PLMs, we aggregate user posts into a single sequence per user (with truncation where necessary) and perform classification at the user level. This design allows PLMs to leverage broader contextual information across a user's writing, while avoiding the limitations of applying classical methods to long-form aggregated text.

Both experiments use a balanced 1:1 ratio of diagnosed to control samples. Our primary goal is to demonstrate the dataset's utility across modeling paradigms rather than to optimize classification performance.

\subsection{Experimental Setup}
\label{subsec:setup}

We frame the task as binary classification: for each mental health condition $C$, the goal is to classify whether a user's data belongs to a user diagnosed with $C$ (positive) or a control user (negative). Control users are randomly sampled from Arabic-speaking Reddit users with no self-reported mental health conditions.

For the classical setup, classification is performed at the post level, where each post inherits the label of its author. For the PLM setup, classification is performed at the user level, where each user is represented by an aggregation of their posts into a single input sequence.

We split the data into 90\% training and 10\% test sets, ensuring no user appears across multiple splits. For post-level experiments, this constraint is enforced by grouping posts by user prior to splitting to prevent leakage.


\subsection{Models}
\label{subsec:models}

\textbf{Classical Classifiers.} We train four traditional machine learning models on TF-IDF bag-of-words features: Naive Bayes (NB), Logistic Regression (LR), Support Vector Machines (SVM), and XGBoost (XGB). These serve as baselines to assess whether simple lexical features can capture condition-specific language patterns.

\textbf{Pre-trained Language Models (PLMs).} We fine-tune two Arabic PLMs:
\begin{itemize}
    \item \textbf{AraBERT-Twitter} (\texttt{aubmindlab/bert-base-arabertv02-twitter}): A BERT model pre-trained on Arabic tweets, well-suited for social media text.
    \item \textbf{CAMeLBERT} (\texttt{CAMeL-Lab/bert-base-arabic-camelbert-mix}): A BERT model pre-trained on a diverse mix of Arabic text including Modern Standard Arabic and dialectal content.
\end{itemize}
For fine-tuning, we use a learning rate of 2e-5, batch size of 16, and train for 3 epochs with early stopping based on validation F1.

\subsection{Evaluation}
\label{subsec:eval}

We evaluate all models using four standard metrics: F1-score, Accuracy, Precision, and Recall. F1-score serves as our primary metric as it balances precision and recall, which is important given the class imbalance across conditions. All metrics are computed on the held-out test set.



% ============================================================
% §7  RESULTS AND ANALYSIS  (~0.75 page)
% STATUS: MISSING : depends on experiments being run.
%
% For a resource paper, the results section should:
%   (1) Present per-condition classification results table
%   (2) Discuss overall trends (which conditions are easier
%       to detect, which are harder)
% ============================================================

\section{Results and Discussion}
\label{sec:results}

Table~\ref{tab:results_summary} summarizes the best F1 scores per condition for both classical classifiers and PLMs. Full per-model breakdowns are provided in Appendix Tables~\ref{tab:results_classical_full} and~\ref{tab:results_plm_full}.

% \input{results_summary.tex}

\subsection{Classical Classifier Performance}

\textbf{Best-performing conditions.} Eating disorder achieves the highest F1 score, with LR reaching .77. Panic disorder also shows strong performance (.74 with SVM). These conditions may exhibit more distinctive linguistic markers in self-disclosure posts.

\textbf{Challenging conditions.} Depression and sleep disorder prove most difficult, with F1 scores ranging from .56--.62. Despite having the largest sample size (1,342 confirmed users for depression), these conditions show lower discriminability, suggesting their linguistic signatures overlap substantially with control users.

\textbf{Model comparison.} Performance differences between classifiers are generally small, with most models achieving similar scores within 2--3 points across conditions.

\subsection{PLM Performance}

CAMeLBERT generally outperforms AraBERT-Twitter across conditions: for anxiety (.70 vs .51), autism (.67 vs .47), bipolar (.62 vs .46), and sleep disorder (.74 vs .54). This suggests CAMeLBERT's diverse pre-training data better captures Reddit's linguistic style and the longer-form nature of Reddit posts. Overall, PLMs provide meaningful gains over classical classifiers for several conditions (e.g., sleep disorder: .62 classical vs .74 PLM), while for others the gap is smaller.


\section{Conclusion}
We introduced CARMA, a large-scale Arabic mental health corpus comprising 3,080 unique users across 16 mental health conditions from Reddit. Our work addresses a critical gap in Arabic NLP resources for mental health research. We developed a dialect-aware extraction pipeline that identifies self-reported diagnoses across Gulf, Egyptian, Levantine, and Modern Standard Arabic, validated through an LLM-as-Judge framework with three independent models achieving moderate inter-annotator agreement.

Our baseline experiments demonstrate that both classical classifiers and pre-trained language models can distinguish diagnosed users from controls, with CAMeLBERT achieving F1 scores up to 0.74 for sleep disorder. The strong performance of CAMeLBERT relative to AraBERT-Twitter highlights the importance of diverse pre-training data for modeling Reddit's longer-form, dialectally varied Arabic content.

We release this corpus to facilitate future research in Arabic computational psychiatry, dialectal NLP, and cross-cultural mental health analysis. We hope CARMA will serve as a foundation for developing more inclusive mental health technologies for the Arabic-speaking world.




\section*{Limitations}
\paragraph{Annotation Pipeline Reliability.}
A small proportion of LLM judge responses failed to parse due to malformed output or API content-safety refusals (Qwen: parse failures on the annotated sample; Jais: up to 16\% raw failure rate, reduced to $<$0.01\% after post-hoc normalization). Unrecoverable responses were treated as abstentions and excluded from the unanimous-agreement corpus.

\paragraph{Sensitive Data Handling.}
This work involves sensitive personal health information derived from social media. Despite our best efforts to anonymize data---including pseudonymizing user identifiers, removing self-disclosure posts, and withholding raw post content---residual re-identification risks may remain. To mitigate misuse, the dataset will be distributed exclusively under a Data Use Agreement (DUA) that restricts access to qualified researchers and prohibits attempts to re-identify individuals.


\section*{Ethical Considerations}
\label{sec:ethics}
This study involves publicly available social media data from users who have voluntarily disclosed their mental health status. We acknowledge the sensitive nature of mental health data and have implemented several safeguards to protect user privacy as listed below.

All data collection adheres to Reddit's Terms of Service and API usage guidelines at the time of collection. We use pre-downloaded public archives from Arctic Shift, which contain posts from publicly accessible Arabic-speaking subreddits, collecting only content that users have chosen to share publicly.

We do not release raw post content or any information that could be used to directly identify individual users. User identifiers are pseudonymized in all released materials. We release only aggregated statistics and model outputs, consistent with best practices in social media mental health research \cite{harrigian2021state}.

The corpus is intended solely for research purposes, specifically to advance Arabic NLP for mental health detection. We recognize that users who disclose mental health conditions on social media may not anticipate their posts being used in research. We mitigate this concern by removing all self-disclosure posts from users' released posting histories, ensuring that the data cannot be used to re-identify individuals based on their diagnostic disclosures. Researchers using this corpus should adhere to ethical guidelines for mental health research and avoid applications that could stigmatize or harm the individuals represented in the data.


% Bibliography entries for the entire Anthology, followed by custom entries
%\bibliography{anthology,custom}
% Custom bibliography entries only
\bibliography{custom}

\appendix

\section{Full Classification Results}
Tables \ref{tab:results_classical_full} and \ref{tab:results_plm_full} show the entire results from our two classification experiments.

\label{app:full_results}

\begin{table}[t]
\centering
\small
\caption{Full post-level classical classifier results on CARMA (Reddit). NB = Naive Bayes, LR = Logistic Regression, SVM = Support Vector Machine, XGB = XGBoost. Best F1 per condition in \textbf{bold}.}
\label{tab:results_classical_full}
\resizebox{\columnwidth}{!}{%
\begin{tabular}{llcccc}
\toprule
\textbf{Condition} & \textbf{Model} & \textbf{F1} & \textbf{Acc.} & \textbf{P} & \textbf{R} \\
\midrule
\multirow{4}{*}{Depression}
 & NB & 0.869 & 0.768 & 0.771 & 0.995 \\
 & LR & 0.858 & 0.754 & 0.774 & 0.963 \\
 & SVM & 0.859 & 0.756 & 0.774 & 0.967 \\
 & XGB & \textbf{0.869} & 0.768 & 0.771 & 0.995 \\
\midrule
\multirow{4}{*}{Anxiety}
 & NB & 0.870 & 0.772 & 0.787 & 0.973 \\
 & LR & 0.852 & 0.747 & 0.789 & 0.926 \\
 & SVM & 0.840 & 0.731 & 0.790 & 0.898 \\
 & XGB & \textbf{0.874} & 0.777 & 0.787 & 0.984 \\
\midrule
\multirow{4}{*}{ADHD}
 & NB & 0.838 & 0.722 & 0.735 & 0.974 \\
 & LR & 0.837 & 0.724 & 0.742 & 0.960 \\
 & SVM & 0.832 & 0.720 & 0.748 & 0.938 \\
 & XGB & \textbf{0.846} & 0.733 & 0.738 & 0.989 \\
\midrule
\multirow{4}{*}{OCD}
 & NB & \textbf{0.811} & 0.683 & 0.691 & 0.983 \\
 & LR & 0.791 & 0.662 & 0.691 & 0.923 \\
 & SVM & 0.758 & 0.630 & 0.693 & 0.837 \\
 & XGB & 0.807 & 0.678 & 0.689 & 0.974 \\
\midrule
\multirow{4}{*}{Sleep dis.}
 & NB & \textbf{0.859} & 0.755 & 0.758 & 0.992 \\
 & LR & 0.850 & 0.743 & 0.759 & 0.967 \\
 & SVM & 0.831 & 0.719 & 0.761 & 0.914 \\
 & XGB & 0.859 & 0.753 & 0.756 & 0.993 \\
\midrule
\multirow{4}{*}{Autism}
 & NB & \textbf{0.854} & 0.748 & 0.757 & 0.981 \\
 & LR & 0.825 & 0.711 & 0.759 & 0.904 \\
 & SVM & 0.773 & 0.653 & 0.763 & 0.783 \\
 & XGB & 0.847 & 0.737 & 0.754 & 0.967 \\
\midrule
\multirow{4}{*}{Panic}
 & NB & \textbf{0.852} & 0.783 & 0.785 & 0.933 \\
 & LR & 0.825 & 0.738 & 0.751 & 0.914 \\
 & SVM & 0.794 & 0.711 & 0.762 & 0.830 \\
 & XGB & 0.818 & 0.717 & 0.721 & 0.946 \\
\midrule
\multirow{4}{*}{Bipolar}
 & NB & \textbf{0.913} & 0.840 & 0.857 & 0.977 \\
 & LR & 0.883 & 0.793 & 0.854 & 0.914 \\
 & SVM & 0.789 & 0.664 & 0.855 & 0.732 \\
 & XGB & 0.900 & 0.819 & 0.852 & 0.954 \\
\midrule
\multirow{4}{*}{PTSD}
 & NB & \textbf{0.889} & 0.801 & 0.801 & 0.999 \\
 & LR & 0.886 & 0.796 & 0.803 & 0.989 \\
 & SVM & 0.836 & 0.726 & 0.802 & 0.873 \\
 & XGB & 0.888 & 0.798 & 0.801 & 0.995 \\
\midrule
\multirow{4}{*}{BPD}
 & NB & 0.681 & 0.550 & 0.534 & 0.941 \\
 & LR & 0.674 & 0.573 & 0.552 & 0.864 \\
 & SVM & 0.636 & 0.574 & 0.565 & 0.730 \\
 & XGB & \textbf{0.685} & 0.582 & 0.557 & 0.888 \\
\midrule
\multirow{4}{*}{Suicidal}
 & NB & \textbf{0.921} & 0.854 & 0.861 & 0.991 \\
 & LR & 0.884 & 0.793 & 0.855 & 0.915 \\
 & SVM & 0.804 & 0.681 & 0.855 & 0.758 \\
 & XGB & 0.887 & 0.799 & 0.857 & 0.920 \\
\midrule
\multirow{4}{*}{Schizophrenia}
 & NB & \textbf{0.611} & 0.441 & 0.441 & 0.990 \\
 & LR & 0.604 & 0.437 & 0.439 & 0.969 \\
 & SVM & 0.581 & 0.453 & 0.440 & 0.856 \\
 & XGB & 0.608 & 0.445 & 0.442 & 0.971 \\
\midrule
\multirow{4}{*}{Eating dis.}
 & NB & \textbf{0.913} & 0.839 & 0.842 & 0.996 \\
 & LR & 0.902 & 0.822 & 0.841 & 0.972 \\
 & SVM & 0.854 & 0.747 & 0.834 & 0.875 \\
 & XGB & 0.891 & 0.803 & 0.836 & 0.953 \\
\bottomrule
\end{tabular}%
}
\end{table}

\input{results_plm.tex}

\begin{table}[h]
\centering
\resizebox{\columnwidth}{!}{%
\begin{tabular}{lc}
\toprule
\textbf{Annotator Pair} & \textbf{CARMA} \\
\midrule
Ann-1 vs.\ Ann-2          & 0.756 \\
\midrule
Ann-1 vs.\ GPT-4.1        & 0.433 \\
Ann-1 vs.\ Qwen3-235B     & 0.563 \\
Ann-1 vs.\ Jais-2-8B      & 0.273 \\
Ann-1 vs.\ LLM Majority   & 0.553 \\
Ann-1 vs.\ LLM Unanimous  & 0.305 \\
\midrule
Ann-2 vs.\ GPT-4.1        & 0.414 \\
Ann-2 vs.\ Qwen3-235B     & 0.541 \\
Ann-2 vs.\ Jais-2-8B      & 0.212 \\
Ann-2 vs.\ LLM Majority   & 0.590 \\
Ann-2 vs.\ LLM Unanimous  & 0.269 \\
\midrule
Mean pairwise $\kappa$      & 0.415 \\
Fleiss' $\kappa$ (all five) & 0.447 \\
\midrule
Overall agreement (\%)      & 67.9\% \\
\bottomrule
\end{tabular}}
\caption{Pairwise Cohen's $\kappa$ between human annotators (Ann-1,
Ann-2) and LLM judges on a stratified 100-post sample from CARMA.
LLM Majority = majority vote across all three models
(reported for comparison only); the corpus retains only posts with
\emph{unanimous} \texttt{<genuine>} agreement across all three
judges. $\kappa$ values follow \citet{landis1977measurement}.}
\label{tab:iaa_human}
\end{table}

%\section{Example Appendix}
%\label{sec:appendix}
%\appendix



%\section{Medication Lexicon}
%\label{app:medications}

%The pipeline detects psychiatric medication mentions using a lexicon of over 300 terms. The lexicon covers:

%\begin{itemize}
%    \item \textbf{Arabic transliterations} with both hamza
%    (\foreignlanguage{arabic}{أ}) and plain alef
%    (\foreignlanguage{arabic}{ا}) variants (e.g.,
%    \foreignlanguage{arabic}{إسيتالوبرام/اسيتالوبرام})
%    \item \textbf{English generic names} (e.g., sertraline,
%    fluoxetine, quetiapine)
%    \item \textbf{Brand names} (e.g., Prozac, Zoloft, Seroquel,
%    Cipralex)
%    \item \textbf{Colloquial terms}
%    (e.g., \foreignlanguage
%    \foreignlanguage{arabic}{مضادات الاكتئاب}
%    [\textit{antidepressants}])
%\end{itemize}

%English terms are matched using word-boundary regex to prevent false positives (e.g., \textit{add} inside \textit{added}), while Arabic terms use substring matching since the \texttt{\textbackslash b} word boundary does not apply to Arabic script.

%\section{Comorbidity Analysis}
%\label{app:comorbidity}

%The pipeline includes a post-processing module that merges all
%per-condition results to identify comorbid presentations. For each
%unique post or comment, the module lists all matched conditions and
%produces three outputs: (1) a unified dataset with a
%\texttt{Conditions} column listing every matched condition per post
%and a \texttt{Condition\_Count} field; (2) an $N \times N$ condition
%co-occurrence matrix showing how many posts matched each pair of
%conditions simultaneously; and (3) a per-condition summary with
%overlap statistics. Users can set a minimum condition threshold
%(e.g., $\geq 2$) to isolate true comorbidity cases.

\section{LLM Classification Prompt}
\label{app:prompt}

We use a two-part prompt structure for LLM-based classification: a \texttt{SYSTEM\_PROMPT} that defines the classifier's role, rules, and output schema, and a \texttt{USER\_TEMPLATE} that instantiates each example with the post, bio, and target condition.

\subsection*{System Prompt}

\begin{tcolorbox}[
  breakable,
  enhanced,
  colback=gray!5,
  colframe=gray!50,
  fontupper=\small\ttfamily,
  coltitle=black,
  attach boxed title to top left={yshift=-2mm, xshift=4mm},
  boxed title style={colback=gray!20, colframe=gray!50}
]
You are a board-certified psychiatrist and clinical researcher specializing in mental health diagnosis. Your native language is Arabic, and you are fluent in English. You understand Arabic dialects, including Gulf, Egyptian, Levantine, and Modern Standard Arabic (MSA), as well as slang, religious expressions, sarcasm, figurative language, and informal social media communication. You are familiar with the cultural and social context of mental health in Arab societies.

Your task is text classification only. Classify based on what the author explicitly states, do not infer or assume a diagnosis.

\medskip\textbf{TASK:}\\
The post was collected using self-disclosure keyword phrases for \{condition\} (\{condition\_ar\}). Determine whether the post is:
\begin{itemize}
  \item \texttt{<genuine>}: a real, first-person self-disclosure of \{condition\}, OR
  \item \texttt{<not\_confirmed>}: NOT a real self-disclosure (false positive).
\end{itemize}

\medskip\textbf{DEFAULT RULE:}\\
When in doubt, label \texttt{<not\_confirmed>}.

\medskip\hrule\medskip

\textbf{REASON TAGS} (list ALL that apply)

\textit{Genuine evidence:}
\begin{itemize}
  \item \texttt{explicit\_self\_disclosure}: Clear first-person statement of having the condition
  \item \texttt{treatment\_or\_medication}: Mentions own medication, therapy, doctor visits, or treatment history
  \item \texttt{symptom\_description\_self}: Describes personal symptoms or lived experience
  \item \texttt{help\_seeking\_self}: Asking for help, resources, or support for their own condition
\end{itemize}

\textit{Not genuine evidence:}
\begin{itemize}
  \item \texttt{metaphor\_or\_exaggeration}: Condition used figuratively or dramatically
  \item \texttt{humor\_or\_sarcasm}: Used jokingly or as a punchline
  \item \texttt{third\_person\_reference}: Talking about someone else's condition
  \item \texttt{negation\_or\_denial}: Explicitly states they do NOT have the condition
  \item \texttt{self\_questioning\_or\_uncertain}: Suspicion without confirmed diagnosis
  \item \texttt{professional\_or\_organizational\_account}: Doctor, clinic, media, or awareness page
  \item \texttt{community\_or\_repost\_account}: Account reposts followers' messages
  \item \texttt{news\_or\_public\_content}: News, educational, or awareness content
\end{itemize}

\medskip\hrule\medskip

\textbf{CLASSIFICATION RULES}

Label \texttt{<genuine>} if the author sincerely states they have \{condition\}, describes their personal experience, seeks help or treatment for their own \{condition\}, or discusses their own medication, symptoms, or diagnosis history.


Label \texttt{<not\_confirmed>} if ANY of the following apply: humor or sarcasm; figurative use; news or third-party reporting; professional or educational posts; community/repost accounts; roleplay or fiction; self-questioning without confirmed diagnosis; explicit negation; third-person discussion; stigmatizing usage; or a vague mention with no sincere self-disclosure intent.

\medskip\hrule\medskip

\textbf{USING THE BIO}

Use BOTH the post and the bio together. A professional, organizational, media, or community bio strongly supports \texttt{<not\_confirmed>}, unless the post contains a clear first-person statement about the author's own diagnosis. A patient identity or direct address to a doctor or support organization supports \texttt{<genuine>}. Specific medications, symptoms, treatment history, or duration strongly support \texttt{<genuine>}.

\medskip\hrule\medskip

\textbf{OUTPUT FORMAT}\\
Output ONLY the following two fields. No preamble, reasoning, or explanation.

\medskip
\texttt{Label: <genuine>} or \texttt{<not\_confirmed>}\\
\texttt{Reasons: one or more tags, comma-separated}
\end{tcolorbox}

\subsection*{User Template}

\begin{tcolorbox}[
  breakable,
  enhanced,
  colback=gray!5,
  colframe=gray!50,
  fontupper=\small\ttfamily,
  coltitle=black,
  attach boxed title to top left={yshift=-2mm, xshift=4mm},
  boxed title style={colback=gray!20, colframe=gray!50}
]
Post: \{post\}\\
Bio: \{bio\}\\
Condition: \{condition\} (\{condition\_ar\})

\medskip
Label:\\
Reasons:
\end{tcolorbox}


\section{Condition Keyword Lists}
\label{app:keywords}

Table~\ref{tab:keywords} lists the Arabic and English condition
keywords used by the pipeline for each of the 16 retained conditions.
Keywords marked with $\dagger$ are \textit{restricted}: they are
permitted only in combination with suffering-class phrases
(\foreignlanguage{arabic}{أعاني من}, \foreignlanguage{arabic}{بعاني من})
to reduce standalone false positives. Orthographic variants (hamza /
hamza-less; taa marbuta \foreignlanguage{arabic}{ة} /
\foreignlanguage{arabic}{ه}) are generated automatically and are not
listed separately.

\begin{table*}[t]
\centering
\small
\caption{Condition keyword inventory (16 conditions retained, 341 keyword forms including English and Arabic variants). $\dagger$ = restricted to suffering phrases only.}
\label{tab:keywords}
\begin{tabular}{p{2.8cm} p{2.2cm} p{10.2cm}}
\toprule
\textbf{Condition} & \textbf{Arabic Label} & \textbf{Representative Keywords (Arabic / English)} \\
\midrule
Schizophrenia &
\foreignlanguage{arabic}{الفصام} &
\foreignlanguage{arabic}{فصام، الفصام، شيزوفرينيا، اضطراب الفصام، انفصام الشخصية، الاضطراب الفصامي الوجداني، سكيزوأفكتيف} /
\textit{schizophrenia, schizoaffective disorder} \\
\addlinespace
Bipolar Disorder &
\foreignlanguage{arabic}{ثنائي القطب} &
\foreignlanguage{arabic}{ثنائي القطب، اضطراب ثنائي القطب، بايبولار، هوس خفيف، مانيا، سيكلوثيميا} /
\textit{bipolar disorder, bipolar I/II, hypomania, mania, cyclothymia} \\
\addlinespace
Depression &
\foreignlanguage{arabic}{الاكتئاب} &
\foreignlanguage{arabic}{اكتئاب، الاكتئاب، اكتئاب حاد $/$ مزمن، اكتئاب ما بعد الولادة، الاكتئاب الموسمي، ديبريشن، سوء المزاج} /
\textit{depression, MDD, dysthymia, postpartum depression, seasonal affective disorder} \\
\addlinespace
Anxiety &
\foreignlanguage{arabic}{اضطراب القلق} &
\foreignlanguage{arabic}{اضطراب القلق، قلق عام $/$ مرضي $/$ مزمن، رهاب اجتماعي، قلق اجتماعي، رهاب الخلاء، قلق الانفصال} /
\textit{generalized anxiety, social anxiety, GAD, agoraphobia, specific phobia} \\
\addlinespace
Panic Disorder &
\foreignlanguage{arabic}{اضطراب الهلع} &
\foreignlanguage{arabic}{نوبة هلع، نوبات هلع، الهلع، اضطراب الهلع، نوبة ذعر} /
\textit{panic attack, panic disorder} \\
\addlinespace
OCD &
\foreignlanguage{arabic}{الوسواس القهري} &
\foreignlanguage{arabic}{وسواس قهري، الوسواس القهري، الوسواس الديني، وسواس الطهارة، ديسمورفيا، اكتناز مرضي} /
\textit{OCD, body dysmorphic disorder, BDD, hoarding disorder, scrupulosity} \\
\addlinespace
PTSD &
\foreignlanguage{arabic}{اضطراب ما بعد الصدمة} &
\foreignlanguage{arabic}{اضطراب ما بعد الصدمة، صدمة نفسية، صدمة معقدة، كرب ما بعد الصدمة} /
\textit{PTSD, C-PTSD, complex PTSD, acute stress disorder} \\
\addlinespace
Eating Disorder &
\foreignlanguage{arabic}{اضطراب الأكل} &
\foreignlanguage{arabic}{فقدان الشهية العصبي، أنوركسيا، الشره المرضي، بوليميا، نهم الطعام، $ARFID$} /
\textit{anorexia, bulimia, binge eating disorder, ARFID, orthorexia} \\
\addlinespace
ADHD &
\foreignlanguage{arabic}{فرط الحركة وتشتت الانتباه} &
\foreignlanguage{arabic}{$ADHD$، فرط حركة، تشتت انتباه، نقص الانتباه، اضطراب الانتباه} /
\textit{ADHD, attention deficit hyperactivity disorder; ADD$\dagger$} \\
\addlinespace
Autism &
\foreignlanguage{arabic}{التوحد} &
\foreignlanguage{arabic}{التوحد، توحد، اضطراب طيف التوحد، طيف التوحد، أسبرجر، متلازمة أسبرجر} /
\textit{autism spectrum disorder, Asperger's} \\
\addlinespace
DID &
\foreignlanguage{arabic}{اضطراب الهوية التفارقي} &
\foreignlanguage{arabic}{اضطراب الهوية التفارقي، تعدد الشخصيات} /
\textit{dissociative identity disorder, DID, multiple personality disorder} \\
\addlinespace
Sleep Disorder &
\foreignlanguage{arabic}{اضطراب النوم} &
\foreignlanguage{arabic}{الأرق، أرق، انقطاع النفس النومي، فرط النوم، الخدار} /
\textit{insomnia, sleep apnea, hypersomnia, narcolepsy} \\
\addlinespace
Trichotillomania &
\foreignlanguage{arabic}{اضطراب شد الشعر} &
\foreignlanguage{arabic}{شد الشعر، نتف الشعر، نتف الجلد، قضم الأظافر} /
\textit{trichotillomania, excoriation disorder, onychophagia} \\
\addlinespace
BPD &
\foreignlanguage{arabic}{اضطراب الشخصية الحدية} &
\foreignlanguage{arabic}{اضطراب الشخصية الحدية، الحدية، بوردرلاين} /
\textit{borderline personality disorder, BPD} \\
\addlinespace
Paranoid PD &
\foreignlanguage{arabic}{اضطراب الشخصية الزوراني} &
\foreignlanguage{arabic}{اضطراب الشخصية الزوراني، الشخصية البارانوية، البارانويا} /
\textit{paranoid personality disorder} \\
\addlinespace
Suicidal Ideation &
\foreignlanguage{arabic}{الأفكار الانتحارية} &
\foreignlanguage{arabic}{أفكار انتحارية، الفكر الانتحاري} /
\textit{suicidal thoughts, suicidal ideation;
\foreignlanguage{arabic}{انتحار}
$\dagger$، \foreignlanguage{arabic}{محاولة انتحار}$\dagger$، \foreignlanguage{arabic}{إيذاء النفس}$\dagger$ / suicide attempt$\dagger$, self-harm$\dagger$} \\
\bottomrule
\end{tabular}
\end{table*}

\section{Self-Disclosure Phrase Inventory}
\label{app:phrases}

Tables~\ref{tab:phrases_universal} and~\ref{tab:phrases_dialect} list
the full self-disclosure phrase inventory (117 universal + 30 dialect-specific not in the universal set = 147 total unique phrases). Universal phrases
(Table~\ref{tab:phrases_universal}) apply across all dialects; they are
organized by semantic type and connector class. Dialect-specific phrases
(Table~\ref{tab:phrases_dialect}) supplement the universal set with
register-appropriate variants for Gulf, Levantine, Egyptian, Maghrebi,
and Modern Standard Arabic (MSA). Orthographic variants (hamza /
hamza-less; taa marbuta \foreignlanguage{arabic}{ة} /
\foreignlanguage{arabic}{ه}) are generated automatically at runtime and
not listed separately.

\begin{table}[t]
\nolinenumbers
\centering
\small
\begin{tabular}{lcc}
\toprule
\multirow{2}{*}{Condition} & \multicolumn{2}{c}{Best F1} \\
\cmidrule(lr){2-3}
 & Classical & PLM \\
\midrule
ADHD & .62 & .64 \\
Anxiety & .61 & .71 \\
Autism & .64 & .67 \\
Bipolar & .62 & .63 \\
BPD & .68 & .55 \\
Depression & .60 & .59 \\
Eating dis. & .78 & --- \\
OCD & .62 & .60 \\
Panic & .74 & .71 \\
PTSD & .64 & .47 \\
Schizophrenia & .65 & --- \\
Sleep dis. & .62 & .74 \\
Suicidal & .64 & .65 \\
\bottomrule
\end{tabular}
\caption{Summary of best F1 scores per condition on CARMA (Reddit). Full results in Appendix Tables~\ref{tab:results_classical_full} and~\ref{tab:results_plm_full}.}
\label{tab:results_summary}
\end{table}
\linenumbers

\begin{table}[t]
\centering
\small
\caption{Universal self-disclosure phrases by type and connector class.
\textbf{B} = \foreignlanguage{arabic}{بـ} connector;
\textbf{M} = \foreignlanguage{arabic}{من} connector;
\textbf{D} = direct (no connector).}
\label{tab:phrases_universal}
\resizebox{\columnwidth}{!}{%
\begin{tabular}{p{2.6cm} c p{5.6cm}}
\toprule
\textbf{Type} & \textbf{Conn.} & \textbf{Phrases} \\
\midrule
Affliction & B &
\foreignlanguage{arabic}{أنا مصاب $/$ انا مصاب، أنا مصابة، أنا مريض، أنا مريضة، كشخص مصاب} \\
\addlinespace
Suffering & M &
\foreignlanguage{arabic}{أعاني من $/$ اعاني من، أنا أعاني من، أنا أتعالج من، أنا أعيش مع، أصبت} \\
\addlinespace
Clinical Diagnosis & B &
\foreignlanguage{arabic}{تم تشخيصي، شخصوني، شخصني، شخصتني، شخصت، شخصني الطبيب، شخصني الدكتور، شخصتني الطبيبة، شخصتني الدكتورة، الأخصائي شخصني، الأخصائية شخصتني} \\
\addlinespace
Possessive & D &
\foreignlanguage{arabic}{عندي، أنا عندي} \\
\addlinespace
Discovery / & B &
\foreignlanguage{arabic}{طلعت مصاب، لقيت نفسي مصاب، اكتشفت إني مصاب، عرفت إن أنا مصاب، اتأكدت إني مصاب، كنت مصاب} \\
Confirmation & D &
\foreignlanguage{arabic}{تشخيصي هو، تشخيصي كان، طلع تشخيصي، طلعت عندي، لقيت نفسي عندي، اكتشفت إن عندي، عرفت إني عندي، اتأكدت إن عندي، بعد التشخيص طلع عندي، التحليل أثبت إن عندي} \\
\addlinespace
Doctor Result & B &
\foreignlanguage{arabic}{أكد الطبيب إصابتي، الدكتور أكد إصابتي، تم التأكيد على إصابتي، تم إثبات إصابتي، الدكتور أكد التشخيص، الطبيب أثبت التشخيص} \\
& D &
\foreignlanguage{arabic}{تشخيص الطبيب هو، تشخيص الدكتور هو، تشخيص حالتي هو، تشخيص المرض هو، تم تشخيص حالتي على أنها، التشخيص الذي حصلت عليه هو، طلع إن التشخيص هو، اتضح إن التشخيص كان} \\
\bottomrule
\end{tabular}
}
\end{table}

\begin{table}[t]
\centering
\small
\caption{Dialect-specific self-disclosure phrases.
\textbf{B} = \foreignlanguage{arabic}{بـ} connector;
\textbf{M} = \foreignlanguage{arabic}{من} connector;
\textbf{D} = direct.}
\label{tab:phrases_dialect}
\resizebox{\columnwidth}{!}{%
\begin{tabular}{p{1.5cm} c p{6.2cm}}
\toprule
\textbf{Dialect} & \textbf{Conn.} & \textbf{Phrases} \\
\midrule
Gulf &
B & \foreignlanguage{arabic}{انا مصاب، مبتلي، مبتلية، ابتليت} \\
& D & \foreignlanguage{arabic}{عندي، فيني، انا متعايش، جاني، جاتني} \\
& M & \foreignlanguage{arabic}{اعاني من} \\
\addlinespace
Levantine &
B & \foreignlanguage{arabic}{انا مصاب} \\
& D & \foreignlanguage{arabic}{عندي، معي، طلعلي، طلعتلي، صار معي، صارت معي، اجاني، اجاتني} \\
& M & \foreignlanguage{arabic}{بعاني من، عم بعاني} \\
\addlinespace
Egyptian &
B & \foreignlanguage{arabic}{انا مصاب، طلعت مصاب، لقيت نفسي مصاب، اكتشفت إني مصاب، اتأكدت إني مصاب} \\
& D & \foreignlanguage{arabic}{عندي، معايا، جالي، جاتلي، طلعلي، طلعت عندي، لقيت نفسي عندي، اكتشفت إن عندي، اتأكدت إن عندي، عرفت إني عندي، أنا بواجه مشكلة اسمها} \\
& M & \foreignlanguage{arabic}{بعاني من، اتشخصت} \\
\addlinespace
Maghrebi &
B & \foreignlanguage{arabic}{انا مصاب} \\
& D & \foreignlanguage{arabic}{عندي، كاين عندي، جاني، جاتني، طلعلي، تشخصت} \\
\addlinespace
MSA &
B & \foreignlanguage{arabic}{أنا مصاب، أنا مصابة، شُخِّصت، شخصت} \\
& D & \foreignlanguage{arabic}{لدي، تشخيصي هو، طلع تشخيصي، أتعايش مع} \\
& M & \foreignlanguage{arabic}{أعاني من، تم تشخيصي} \\
\bottomrule
\end{tabular}
}
\end{table}

\section{Connector Logic and Orthographic Variants}
\label{app:connectors}

Table~\ref{tab:connectors} summarizes the three connector types used by
the pipeline, showing how each phrase class links to the condition
keyword and how the definite article \foreignlanguage{arabic}{ال}
triggers the \foreignlanguage{arabic}{بال} fusion rule.

\begin{table}[h]
\centering
\small
\caption{Connector rules: phrase type, connector morpheme, keyword form generated, and example.}
\label{tab:connectors}
\resizebox{\columnwidth}{!}{%
\begin{tabular}{p{2.0cm} c p{3.0cm} p{3.5cm}}
\toprule
\textbf{Phrase Type} & \textbf{Conn.} & \textbf{Keyword Form} & \textbf{Example} \\
\midrule
Affliction, Diagnosis, Confirmation &
\foreignlanguage{arabic}{بـ} &
Base kw: \foreignlanguage{arabic}{بـ}+kw \newline
\foreignlanguage{arabic}{ال}-kw: \foreignlanguage{arabic}{بال}+stem \newline
English: \foreignlanguage{arabic}{بـ} kw &
\foreignlanguage{arabic}{مصاب باكتئاب} \newline
\foreignlanguage{arabic}{مصاب بالاكتئاب} \newline
\foreignlanguage{arabic}{مصاب بـ $PTSD$} \\
\addlinespace
Suffering &
\foreignlanguage{arabic}{من} &
Phrase ends in \foreignlanguage{arabic}{من}: bare kw \newline
Phrase lacks \foreignlanguage{arabic}{من}: \foreignlanguage{arabic}{من}+kw &
\foreignlanguage{arabic}{أعاني من الاكتئاب} \newline
\foreignlanguage{arabic}{عم بعاني من اكتئاب} \\
\addlinespace
Possessive, Discovery &
Direct &
Bare or definite kw &
\foreignlanguage{arabic}{عندي اكتئاب} \newline
\foreignlanguage{arabic}{تشخيصي هو الفصام} \\
\bottomrule
\end{tabular}
}
\end{table}

\noindent The pipeline generates all orthographic variants automatically.
Table~\ref{tab:ortho} lists the variant pairs handled.

\begin{table}[h]
\centering
\small
\caption{Orthographic variant pairs generated automatically.}
\label{tab:ortho}
\begin{tabular}{lll}
\toprule
\textbf{Feature} & \textbf{Variant A} & \textbf{Variant B} \\
\midrule
Hamza on alef & \foreignlanguage{arabic}{أنا} & \foreignlanguage{arabic}{انا} \\
Hamza in conjunction & \foreignlanguage{arabic}{إن} & \foreignlanguage{arabic}{ان} \\
Taa marbuta & \foreignlanguage{arabic}{مصابة} & \foreignlanguage{arabic}{مصابه} \\
Taa marbuta & \foreignlanguage{arabic}{مريضة} & \foreignlanguage{arabic}{مريضه} \\
Hamza on \foreignlanguage{arabic}{ا} & \foreignlanguage{arabic}{أعاني} & \foreignlanguage{arabic}{اعاني} \\
\bottomrule
\end{tabular}
\end{table}


\section{Agreement Computation Details}
\label{app:agreement_details}

\subsection*{Sample sizes and parse failures}

All pairwise $\kappa$ values reported in Tables~\ref{tab:iaa}
and~\ref{tab:iaa_human} are computed on the subset of posts where both
raters returned a valid label. Human-to-human agreement always uses the
full $n=100$. Jais-2-8B-Chat produced parse failures on 16\% of CARMA
posts; any comparison involving Jais-2-8B excludes those posts,
yielding $n=84$ for affected pairs. GPT-4.1 produced no parse failures
on CARMA; Qwen3-235B produced no parse failures. Ensemble metrics (LLM
Majority, LLM Unanimous) and multi-rater Fleiss' $\kappa$ are computed
on complete cases only, posts where all three LLM judges returned a
valid label (CARMA $n=84$), as both metrics require
a label from every judge. Table~\ref{tab:sample_sizes} summarises the
effective $n$ for each annotator pair.

\begin{table}[h]
\centering
\small
\begin{tabular}{lc}
\toprule
\textbf{Annotator Pair} & \textbf{CARMA} \\
\midrule
Ann-1 vs.\ Ann-2          & 100 \\
\midrule
Ann-1/2 vs.\ GPT-4.1      & 100 \\
Ann-1/2 vs.\ Qwen3-235B   & 100 \\
Ann-1/2 vs.\ Jais-2-8B    &  84 \\
Ann-1/2 vs.\ LLM Majority &  84 \\
Ann-1/2 vs.\ LLM Unanimous&  84 \\
\midrule
GPT-4.1 vs.\ Qwen3-235B   & 100 \\
GPT-4.1 vs.\ Jais-2-8B    &  84 \\
Qwen3-235B vs.\ Jais-2-8B &  84 \\
\midrule
Fleiss' $\kappa$ (all)     &  84 \\
\bottomrule
\end{tabular}
\caption{Effective sample size $n$ for each annotator pair after
excluding posts with parse failures from the relevant model(s).}
\label{tab:sample_sizes}
\end{table}

\subsection*{Kappa interpretation scale}

Cohen's $\kappa$ and Fleiss' $\kappa$ values are interpreted following
\citet{landis1977measurement}: $\kappa \leq 0.20$ slight,
$0.21$--$0.40$ fair, $0.41$--$0.60$ moderate, $0.61$--$0.80$
substantial, $0.81$--$1.00$ almost perfect agreement.

\subsection*{Ensemble label definitions}

\textbf{LLM Majority} assigns \texttt{<genuine>} to a post if at least
two of the three LLM judges return \texttt{<genuine>}. This metric is
reported in Table~\ref{tab:iaa_human} for comparison purposes only.
\textbf{LLM Unanimous} assigns \texttt{<genuine>} only when all three
judges agree. The final CARMA corpus retains exclusively
posts that satisfy the unanimous criterion, prioritising precision over
recall to minimise false positives.


\end{document}
